%\documentclass{...}

\input{./def/yf-formatting}

\usepackage{amsfonts, amsmath, amssymb, amsthm}
\usepackage{comment}
\usepackage{graphicx}
\usepackage{ifthen}
\usepackage{latexsym}
%\usepackage{times}
\usepackage[normalem]{ulem}

\input{./def/yf-def}

\def\dom{\prec}
\def\T{\mathcal{T}}

\newboolean{solver}\setboolean{solver}{true}
%\newboolean{solver}\setboolean{solver}{false}
\ifthenelse{\boolean{solver}}{\includecomment{sol}}{\excludecomment{sol}}

\begin{document}

\section*{CE7456: Exercise List 3}
Prepared by Yufei Tao \\

% \begin{center}
%     \uline{Classroom Discussion}
% \end{center}

\boxminipg{0.95\linewidth}{
    These exercises are designed to help you practice the techniques taught in class. Attempt them ``at your own risk'' because they are research-oriented and can be rather challenging. Solutions will not be provided in written form. I will be happy to act as a collaborator and sounding board --- I look forward to discussing your creative approaches, refining your key steps, and walking through your attempted solutions with you.
}

\extraspacing {\bf Problem 1.} In the DDS (deferred data structuring) lecture, we required each query answer to be stored in $O(1)$ words. This problem lets you explore how to extend the DDS technique to handle reporting queries.

\vgap

Let $S$ be a set of $n$ integers. Given an interval $q$, a {\em range reporting query} returns $S \cap q$. Describe an algorithm that can answer any $r \le n$ such queries in $O(n \log r + \sum_{i=1}^r k_i)$ time, where $k_i$ is the number of elements reported by the $i$-th query.

\extraspacing {\bf Problem 2.} Recall that, in the IQS version of range reporting, the input is a set $S$ of $n$ integers, each of which is associated with a positive integer weight. Given an interval $q$ and an integer $s$, a query returns $s$ independent weighted samples from $S \cap q$. Describe an algorithm that can answer any $r \le n$ such queries in $O(n \log r + \sum_{i=1}^r s_i)$ time, where $s_i$ is the number of samples demanded by the $i$-th query.

\vgap

(Hint: First consider $s_i = 1$ for all $i \in [r]$ and argue that the problem is $(\log n, \log n)$-spectrum indexable. Then, extend your solution to allow arbitrary $s_1, s_2, ..., s_r$.)

\extraspacing {\bf Problem 3.} Consider a $(\log n, \log n)$-spectrum indexable problem. Prove: If we can support any $r \le \sqrt{n}$ queries in $O(n \log r)$ time, then we can support any $r \le n$ queries in $O(n \log r)$ time.

\extraspacing {\bf Problem 4*.} Consider a $(\log n, \sqrt{n})$-spectrum indexable problem. Prove: we can support any $r \le \sqrt{n} \log n$ queries in $O(n \log r)$ time.

\extraspacing {\bf Problem 5.} Let $S$ be a set of $n$ integers stored in a binary search tree (BST) $\T$ where each node stores a pointer to its parent. Explain how to use $\T$ to find the predecessor of an arbitrary integer $q$ in $O(\log \Delta)$ time, where $\Delta$ is the number of elements in $S$ that are no greater than $q$.

\extraspacing {\bf Problem 6.} Let us modify the structure discussed in the ``learned index'' lecture by always fixing $a = 1$ and $b = U$. Prove: Theorem 1 in the lecture notes still holds on the modified structure.

\end{document}
